% 从一阶基础不等式到临界、Morrey 和 Rellich；保留旧交叉引用。
\Needspace{12\baselineskip}
\section{Sobolev 嵌入：次临界、临界、连续性与紧性的详细证明}\label{sec:sobolev-embedding}

\subsection{先区分包含、连续嵌入与紧嵌入}

$X\subset Y$ 只说明每个 $X$ 元素也属于 $Y$。$X\hookrightarrow Y$ 还要求恒等映射连续，即存在与 $u$ 无关的 $C$，使 $\|u\|_Y\le C\|u\|_X$。这保证 $X$ 中强收敛传递到 $Y$。$X\Subset Y$ 表示紧嵌入：$X$ 中每个有界序列都有 $Y$ 中的强收敛子列。

接下来的证明按四条路线展开：$p<d$ 得更高可积性，$p=d$ 得任意有限次幂可积性，$p>d$ 得 H\"older 连续代表元，最后以平移估计证明紧性。需要一般边界延拓时，完整证明统一采用上一节已经构造过的有界 $C^1$ 区域。

\begin{theorem}[次临界一阶 Sobolev 嵌入]\label{thm:sobolev-main}
设整数 $d\ge2$，$1\le p<d$，$p^*=dp/(d-p)$。对任意开集 $\Omega\subset\R^d$，
\begin{equation}
 \boxed{\|u\|_{p^*}\le C_{d,p}\|\nabla u\|_p,
       \qquad u\in W^{1,p}_0(\Omega).}
 \label{eq:sobolev-main}
\end{equation}
在全空间 $W^{1,p}(\R^d)$ 上也有同一估计。若 $\Omega$ 有界且为 $C^1$，一般非零边界函数满足
\begin{equation}
 \|u\|_{p^*}\le C_{\Omega,p}\big(\|u\|_p+\|\nabla u\|_p\big),
 \qquad u\in W^{1,p}(\Omega).
 \label{eq:sobolev-nonzero}
\end{equation}
\end{theorem}

\subsection{基础积分引理：将不同坐标缺失的函数乘起来}

\begin{lemma}[坐标乘积积分估计]\label{lem:sobolev-coordinate}
对每个 $i=1,\ldots,d$，设非负可积函数 $A_i(\widehat x_i)$ 只依赖除 $x_i$ 以外的 $d-1$ 个变量，则
\begin{equation}
 \int_{\R^d}\prod_{i=1}^d A_i(\widehat x_i)^{1/(d-1)}\dd x
 \le\prod_{i=1}^d
       \left(\int_{\R^{d-1}}A_i\dd\widehat x_i\right)^{1/(d-1)}.
 \label{eq:sobolev-coordinate}
\end{equation}
\end{lemma}
\begin{proof}
对维数归纳。$d=2$ 时是 $\iint A_1(x_2)A_2(x_1)=\int A_1\int A_2$。

设 $d\ge3$，并假设 $d-1$ 维成立。先沿 $x_d$ 积分，$A_d$ 与它无关。对其余 $d-1$ 个因子用指数全为 $d-1$ 的 H\"older，令
\[
 G_i(\widehat x_i')=\int_\R A_i(\widehat x_i)\dd x_d
 \quad(1\le i<d),\qquad x'=(x_1,\ldots,x_{d-1}).
\]
则原积分不超过
\[
 \int_{\R^{d-1}} A_d(x')^{1/(d-1)}
                \prod_{i<d}G_i(\widehat x_i')^{1/(d-1)}\dd x'.
\]
设 $H(x')=\prod_{i<d}G_i^{1/(d-2)}$，上式的后半乘积为 $H^{(d-2)/(d-1)}$。再用共轭指数 $d-1$、$(d-1)/(d-2)$ 的 H\"older，上式不超过
\[
 \left(\int A_d\right)^{1/(d-1)}
 \left(\int H\right)^{(d-2)/(d-1)}.
\]
对 $H$ 使用 $d-1$ 维归纳假设，即 $\int H\le\prod_{i<d}(\int G_i)^{1/(d-2)}$。而 Tonelli 给出 $\int G_i=\int A_i$。代入即可得到所需的全部 $d$ 个因子。
\end{proof}

\subsection{第一步：从微积分基本定理证明 \texorpdfstring{$W^{1,1}$}{W11} 估计}

设 $v\in C_c^1(\R^d)$。固定除 $x_i$ 以外的变量，从负无穷积分，因紧支撑有
\[
 |v(x)|=\left|\int_{-\infty}^{x_i}\partial_i v(x_1,\ldots,s,\ldots,x_d)\dd s\right|
 \le A_i(\widehat x_i):=\int_\R|\partial_i v|\dd x_i.
\]
将 $d$ 个不等式相乘并取 $1/(d-1)$ 次幂，得
\[
 |v(x)|^{d/(d-1)}\le\prod_{i=1}^d A_i(\widehat x_i)^{1/(d-1)}.
\]
积分并用刚证明的引理，再取 $(d-1)/d$ 次幂：
\begin{align*}
 \|v\|_{d/(d-1)}
 &\le\prod_{i=1}^d\|\partial_i v\|_1^{1/d}\\
 &\le\frac1d\sum_{i=1}^d\|\partial_i v\|_1
 \le C_d\|\nabla v\|_1.
\end{align*}
于是
\begin{equation}
 \boxed{\|v\|_{d/(d-1)}\le C_d\|\nabla v\|_1.}
 \label{eq:sobolev-w11}
\end{equation}
这是整个次临界嵌入证明的起点，且已经覆盖 $p=1$。

\textbf{二维单独核对。}此时 $|v(x,y)|$ 同时不超过 $A(y)$ 与 $B(x)$，所以
\[
 |v|^2\le A(y)B(x),\qquad
 \|v\|_2^2\le\|\partial_1v\|_1\|\partial_2v\|_1.
\]

\textbf{三维单独核对。}令 $A(y,z)=\int|\partial_1v|\dd x$、$B(x,z)=\int|\partial_2v|\dd y$、$C(x,y)=\int|\partial_3v|\dd z$。则
\[
 \iiint|v|^{3/2}\le\iiint(ABC)^{1/2}
 \le(\|\partial_1v\|_1\|\partial_2v\|_1\|\partial_3v\|_1)^{1/2}.
\]
最后一个估计就是引理在 $d=3$ 时的三次 Cauchy--Schwarz。取 $2/3$ 次幂得到 $\|v\|_{3/2}\le C\|\nabla v\|_1$。

\subsection{第二步：为什么必须选这个幂指数}

设 $1<p<d$，先取 $u\in C_c^\infty(\R^d)$。希望把基础估计应用于 $v=|u|^\gamma$。左边会出现 $\|u\|_{\gamma d/(d-1)}^\gamma$；右边的导数乘积在 H\"older 后会出现 $\|u\|_{(\gamma-1)p'}^{\gamma-1}$。为了两侧使用同一个函数范数，要求
\[
 \frac{\gamma d}{d-1}=(\gamma-1)p'.
\]
这不是猜测，而是解一个代数方程。代入 $p'=p/(p-1)$，交叉相乘得
\[
 \gamma d(p-1)=p(d-1)(\gamma-1),\qquad
 \gamma(d-p)=p(d-1).
\]
所以
\begin{equation}
 \gamma=\frac{p(d-1)}{d-p}>1,\qquad
 \frac{\gamma d}{d-1}=(\gamma-1)p'=p^*=\frac{dp}{d-p}.
 \label{eq:sobolev-power-choice}
\end{equation}
因为 $\gamma>1$，$|s|^\gamma$ 在零点也 $C^1$，故 $v\in C_c^1$；不需要宣称它无穷可微。逐项计算得
\begin{align*}
 \|u\|_{p^*}^{\gamma}
 &\le C_d\|\nabla(|u|^\gamma)\|_1\\
 &=C_d\gamma\int|u|^{\gamma-1}|\nabla u|\dd x\\
 &\le C_d\gamma\|\nabla u\|_p
      \left(\int|u|^{(\gamma-1)p'}\dd x\right)^{1/p'}\\
 &=C_d\gamma\|\nabla u\|_p\|u\|_{p^*}^{\gamma-1}.
\end{align*}
若 $u$ 非零可除去 $\|u\|_{p^*}^{\gamma-1}$；若为零结论直接成立。这就证明了光滑函数版本。三维 $p=2$ 对应 $\gamma=4$：
\[
 \|u\|_6^4\le C\int|u|^3|\nabla u|
 \le C\|u\|_6^3\|\nabla u\|_2.
\]

\subsection{第三步：光滑函数的估计怎样推广到 Sobolev 函数}

\textbf{零边界情形。}若 $u\in W^{1,p}_0(\Omega)$，按定义取 $u_n\in C_c^\infty(\Omega)$，使其在 $W^{1,p}$ 范数中趋于 $u$。将各 $u_n$ 零延拓，已证估计用于差值：
\[
 \|u_n-u_m\|_{p^*}\le C\|\nabla u_n-\nabla u_m\|_p\to0.
\]
由 $L^{p^*}$ 完备性得到某个 $w\in L^{p^*}$。还须确认 $w=u$：对任意有限测度 $K\Subset\Omega$，
\[
 \|u-w\|_{L^1(K)}
 \le |K|^{1-1/p}\|u-u_n\|_{L^p(K)}
     +|K|^{1-1/p^*}\|u_n-w\|_{L^{p^*}(K)}\to0.
\]
故 $u=w$ a.e.，用可数穷竭覆盖区域即可。最后范数收敛允许在不等式两侧取极限，证明~\eqref{eq:sobolev-main}。

\textbf{全空间情形。}上一节已经构造了 $C_c^\infty(\R^d)$ 的 $W^{1,p}$ 逼近，因此逐字重复上面的 Cauchy 与极限识别论证即可。

\textbf{一般边界情形。}对有界 $C^1$ 区域，使用定理~\ref{thm:weak-extension} 的延拓 $Eu$，则
\[
 \|u\|_{L^{p^*}(\Omega)}
 \le\|Eu\|_{L^{p^*}(\R^d)}
 \le C\|\nabla Eu\|_{L^p(\R^d)}
 \le C_\Omega\|u\|_{W^{1,p}(\Omega)}.
\]
延拓中的截断产生零阶项，所以不能删掉 $\|u\|_p$。非零常数也是立即可见的反例。至此定理~\ref{thm:sobolev-main} 的三个版本全部证明。

\subsection{缩放推导临界指数，并构造不能改进的反例}

令 $u_\lambda(x)=u(\lambda x)$，变量替换给出
\[
 \|u_\lambda\|_q=\lambda^{-d/q}\|u\|_q,\qquad
 \|\nabla u_\lambda\|_p=\lambda^{1-d/p}\|\nabla u\|_p.
\]
若全空间的纯梯度估计 $\|u\|_q\le C\|\nabla u\|_p$ 对所有尺度成立且 $C$ 不变，则必须 $-d/q=1-d/p$，即 $q=p^*$。

在有界区域内部一点 $x_0$ 附近，取非零 $\phi\in C_c^\infty(B_1)$，定义
\[
 u_n(x)=n^{d/p-1}\phi(n(x-x_0)),
\]
对充分大的 $n$ 支撑落在区域内。计算
\begin{equation}
 \|\nabla u_n\|_p=\|\nabla\phi\|_p,\quad
 \|u_n\|_p=n^{-1}\|\phi\|_p,\quad
 \|u_n\|_q=n^{d/p-1-d/q}\|\phi\|_q.
 \label{eq:sobolev-concentration}
\end{equation}
若 $q>p^*$，最后指数为正，$L^q$ 范数发散，所以这种更高指数的连续嵌入一般不成立。若 $q=p^*$，范数恒为正，而 $u_n\to0$ a.e.，故不可能有 $L^{p^*}$ 强收敛子列。这同时解释端点连续但一般不紧。

\subsection{临界情况 \texorpdfstring{$p=d$}{p=d}：有限次幂可以，端点无穷不可以}

\begin{theorem}[临界有限指数嵌入]\label{thm:sobolev-critical}
若 $d\ge2$，$\Omega$ 有界 $C^1$，则对每个 $1\le q<\infty$，
\[
 \|u\|_q\le C_{\Omega,q}\|u\|_{W^{1,d}},
 \qquad u\in W^{1,d}(\Omega).
\]
\end{theorem}
\begin{proof}
$q\le d$ 时由有限测度上的 $L^d\hookrightarrow L^q$ 直接成立。若 $q>d$，选
\[
 r=\frac{dq}{d+q},\qquad1<r<d,\qquad\frac{dr}{d-r}=q.
\]
有限测度上的 H\"older 给出 $W^{1,d}\hookrightarrow W^{1,r}$。对 $r$ 使用已证的次临界一般边界嵌入：
\[
 \|u\|_q\le C_{\Omega,r}\|u\|_{W^{1,r}}
 \le C_{\Omega,q}\|u\|_{W^{1,d}}.
\]
这说明每个有限 $q$ 都可以，但常数依赖 $q$，没有推出 $q=\infty$ 的估计。
\end{proof}

\begin{example}[临界无界函数的完整检查]\label{ex:sobolev-loglog}
在 $d\ge2$、原点附近取 $u(x)=\log\log(e/|x|)$，再乘支撑于 $B_1$、原点附近为 1 的光滑截断。令 $r=|x|$，有
\[
 |u'(r)|=\frac1{r\log(e/r)},\qquad
 \int_0^\delta|u'(r)|^dr^{d-1}\dd r
 =\int_0^\delta\frac{\dd r}{r\log^d(e/r)}<\infty.
\]
最后的积分以 $s=\log(e/r)$ 化为 $\int^\infty s^{-d}\dd s$，因 $d>1$ 收敛。函数本身的 $d$ 次幂积分也收敛：同一换元后被指数衰减 $e^{-ds}$ 控制。去掉 $B_\varepsilon$ 分部积分时，内边界项不超过 $C\varepsilon^{d-1}\log\log(e/\varepsilon)\to0$，所以候选导数确实是弱导数，无额外原点质量。因此 $u\in W^{1,d}_0(B_1)$，但 $u$ 本质无界。特别二维 $H^1$ 不嵌入 $L^\infty$。
\end{example}

\subsection{Morrey 证明的核心：从线段平均得到奇核估计}

设 $u\in C_c^\infty(\R^d)$，$B_r(x)$ 为球，记 $u_B=|B|^{-1}\int_Bu$。从沿线段的基本定理出发，
\[
 |u(x)-u(x+s\omega)|\le\int_0^s|\nabla u(x+t\omega)|\dd t.
\]
在球上取平均，极坐标和 Tonelli 给出
\begin{align*}
 |u(x)-u_{B_r(x)}|
 &\le\frac1{|B_r|}\int_{\mathbb S^{d-1}}\int_0^r
       \int_0^s|\nabla u(x+t\omega)|\dd t\,s^{d-1}\dd s\dd\omega\\
 &\le C_d\int_{\mathbb S^{d-1}}\int_0^r
                   |\nabla u(x+t\omega)|\dd t\dd\omega\\
 &=C_d\int_{B_r(x)}\frac{|\nabla u(z)|}{|x-z|^{d-1}}\dd z.
\end{align*}
第二行使用 $|B_r|^{-1}\int_t^rs^{d-1}\dd s\le C_d$。这个核的分母次数来自极坐标体积元 $t^{d-1}$，不能凭空省略。

若 $p>d$，则 $p'=p/(p-1)$ 满足 $(d-1)p'<d$。H\"older 的核积分为
\begin{align*}
 \left(\int_{B_r(x)}|x-z|^{-(d-1)p'}\dd z\right)^{1/p'}
 &=\left(|\mathbb S^{d-1}|\int_0^r
                 t^{d-1-(d-1)p'}\dd t\right)^{1/p'}\\
 &=C_{d,p}r^{1-d/p}.
\end{align*}
因此
\begin{equation}
 |u(x)-u_{B_r(x)}|\le C_{d,p}r^{1-d/p}\|\nabla u\|_{L^p(B_r(x))}.
 \label{eq:sobolev-morrey-average}
\end{equation}
这里恰好需要 $p>d$；当 $p=d$ 核积分临界发散，这说明证明在哪里失效。

\subsection{Morrey 连续嵌入：两点估计、上界与代表元}

\begin{theorem}[Morrey 不等式]\label{thm:sobolev-morrey}
若 $d<p<\infty$，$\alpha=1-d/p$，则 $W^{1,p}(\R^d)$ 的每个 a.e. 等价类有唯一连续代表元，满足
\begin{equation}
 [u]_{C^{0,\alpha}}:=\sup_{x\ne y}\frac{|u(x)-u(y)|}{|x-y|^\alpha}
 \le C\|\nabla u\|_p,\qquad
 \|u\|_\infty\le C(\|u\|_p+\|\nabla u\|_p).
 \label{eq:sobolev-morrey}
\end{equation}
有界 $C^1$ 区域上有 $W^{1,p}(\Omega)\hookrightarrow C^{0,\alpha}(\overline\Omega)$，估计由完整 Sobolev 范数控制。
\end{theorem}
\begin{proof}
\textbf{两点之差。}先设 $u$ 光滑。取 $x\ne y$、$r=2|x-y|$，令 $E=B_r(x)\cap B_r(y)$。交集包含半径 $3r/4$、中心为 $(x+y)/2$ 的球，因此 $|E|\ge c_dr^d$。对 $z\in E$ 写
\[
 |u(x)-u(y)|\le |u(x)-u(z)|+|u(y)-u(z)|,
\]
在 $E$ 上平均。第一项不超过 $C_dr^{-d}\int_{B_r(x)}|u(x)-u(z)|\dd z$，第二项类似。沿上一小节的线段平均计算再用 H\"older，即得
$|u(x)-u(y)|\le C r^\alpha\|\nabla u\|_p=C'|x-y|^\alpha\|\nabla u\|_p$。

\textbf{本身的上界。}取半径 1 的球，利用~\eqref{eq:sobolev-morrey-average} 及 $|u_{B_1(x)}|\le |B_1|^{-1/p}\|u\|_p$，得到一致的 $\|u\|_\infty$ 估计。

\textbf{推广与代表元。}取 $u_n\in C_c^\infty$ 使其在 $W^{1,p}$ 中趋于 $u$。上述两个估计用于 $u_n-u_m$，说明它们在完整 $C^{0,\alpha}$ 范数下 Cauchy，特别一致收敛到某个连续函数 $\widetilde u$。它在局部 $L^p$ 的极限与 $u$ 一致，所以是代表元；不等式取极限保留。两个连续代表元若 a.e. 相同，差若某点非零则在一个球内非零，与 a.e. 相同矛盾，故代表元唯一。一般边界情形先延拓到全空间，再限制至闭包即可。
\end{proof}

当 $p=\infty$，全空间 $W^{1,\infty}$ 具有 Lipschitz 代表元。可在光滑卷积上用线段基本定理得到 $|u_\varepsilon(x)-u_\varepsilon(y)|\le\|\nabla u\|_\infty|x-y|$，再用局部一致紧性与局部 $L^1$ 收敛识别代表元；这里无需 $W^{1,\infty}$ 范数的光滑稠密性。凸区域内也能直接沿线段估计；一般区域应考虑延拓或域的几何性质。

\subsection{Poincar\'e：零边界版本的独立证明}

\begin{theorem}[零边界 Poincar\'e]\label{thm:sobolev-poincare}
若 $\Omega$ 有界，$1\le p<\infty$，则
\begin{equation}
 \|u\|_p\le C_\Omega\|\nabla u\|_p,
 \qquad u\in W^{1,p}_0(\Omega).
 \label{eq:sobolev-poincare}
\end{equation}
\end{theorem}
\begin{proof}
将区域平移后放进边长 $L$ 的立方体 $Q=(0,L)^d$。对 $u\in C_c^\infty(\Omega)$ 零延拓，有
\[
 u(x_1,x')=\int_0^{x_1}\partial_1u(s,x')\dd s.
\]
H\"older（$p=1$ 直接积分）给出
\[
 |u(x_1,x')|^p\le L^{p-1}\int_0^L|\partial_1u(s,x')|^p\dd s.
\]
再积分 $x_1,x'$ 得 $\|u\|_p^p\le L^p\|\partial_1u\|_p^p$。按 $W^{1,p}_0$ 定义逼近并取极限即可。这个证明不需要边界光滑，也不依赖后面的紧性定理。
\end{proof}

因此 $\|\nabla u\|_p$ 在零边界空间上是等价范数；在一般 $W^{1,p}$ 上仍只是半范数。三维的 Sobolev 嵌入再结合有限测度的 $L^6\hookrightarrow L^2$ 也能推出 $p=2$ 的零边界 Poincar\'e，但沿线积分证明适用范围更直接。

% 延续嵌入章节：空间紧性、零均值 Poincare 和高阶补充。
\subsection{Rellich 的第一步：一致平移估计与一致光滑逼近}

对 $f\in W^{1,p}(\R^d)$，$1\le p<\infty$，沿向量 $h$ 的线段积分并用 Jensen，得到
\begin{equation}
 \|f(\cdot-h)-f\|_p\le |h|\|\nabla f\|_p.
 \label{eq:sobolev-translation-uniform}
\end{equation}
光滑时逐点有 $f(x-h)-f(x)=-\int_0^1\nabla f(x-th)\cdot h\dd t$，积分即可；一般情况用光滑稠密性推广。将此式用于卷积平均，得到定量估计
\begin{equation}
 \|f*\rho_\varepsilon-f\|_p
 \le C_\rho\varepsilon\|\nabla f\|_p.
 \label{eq:sobolev-uniform-mollify}
\end{equation}
与普通近似恒等不同，这个误差对\textbf{整个梯度范数有界的函数族一致趋零}，它是紧性证明的核心。

设 $u_n$ 在有界 $C^1$ 区域的 $W^{1,p}$ 中一致有界，用固定的有界支撑延拓得到 $f_n=Eu_n$，则 $\|f_n\|_{W^{1,p}}\le M$，所有支撑都在一个固定紧集内。令 $f_{n,\varepsilon}=\rho_\varepsilon*f_n$。对固定 $\varepsilon>0$，由 H\"older，
\[
 \|f_{n,\varepsilon}\|_\infty\le\|\rho_\varepsilon\|_{p'}M,
 \quad
 \|\nabla f_{n,\varepsilon}\|_\infty\le\|\nabla\rho_\varepsilon\|_{p'}M.
\]
这些函数支撑统一有界，且一致有界、一致 Lipschitz，因此在一致范数下可抽出收敛子列。

为看清最后一步，可用有限网格证明：给定误差 $\eta$，用有限个半径足够小的球覆盖固定支撑集，Lipschitz 界保证每个球内振荡小于 $\eta$；中心点上的函数值位于固定有界区间，将这些值以步长 $\eta$ 量化，只有有限种组合。每种组合给出一个一致范数的有限精度近似。因此该函数族完全有界，再用连续函数空间完备性即可抽取一致收敛子列。这就是此处所需的 Arzel\`a--Ascoli 论证。

\Needspace{11\baselineskip}
\subsection{Rellich 的第二步：对角子列与原序列的强收敛}

\begin{theorem}[Rellich--Kondrachov 紧嵌入]\label{thm:sobolev-rellich}
设 $\Omega$ 有界 $C^1$，$1\le p<\infty$。则
\begin{enumerate}
 \item $W^{1,p}(\Omega)\Subset L^p(\Omega)$；
 \item 若 $p<d$，则 $W^{1,p}(\Omega)\Subset L^q(\Omega)$ 对每个 $1\le q<p^*$ 成立；
 \item 若 $p=d\ge2$，则对每个有限 $q$ 紧嵌入 $L^q$；
 \item 若 $p>d$，则紧嵌入 $C^{0,\beta}(\overline\Omega)$，$0<\beta<1-d/p$。
\end{enumerate}
\end{theorem}
\begin{proof}
\textbf{先证明 $L^p$ 紧性。}取 $\varepsilon_j\downarrow0$。对 $j=1$ 抽出使 $f_{n,\varepsilon_1}$ 一致收敛的子列；再在其中抽出使 $f_{n,\varepsilon_2}$ 一致收敛的子列，依次进行。取对角子列后，对每个固定 $j$，相应光滑函数都一致收敛，从而也在固定有界支撑集上 $L^p$ 收敛。

对角子列仍记 $f_n$。给定 $\eta>0$，先选 $j$ 使 $2C_\rho M\varepsilon_j<\eta/2$，再令 $m,n$ 足够大，使光滑项的 $L^p$ 距离小于 $\eta/2$。于是
\begin{align*}
 \|f_n-f_m\|_p
 &\le\|f_n-f_{n,\varepsilon_j}\|_p
     +\|f_{n,\varepsilon_j}-f_{m,\varepsilon_j}\|_p
     +\|f_{m,\varepsilon_j}-f_m\|_p\\
 &<\eta.
\end{align*}
故原函数子列在 $L^p$ 中 Cauchy，由完备性强收敛；限制回 $\Omega$ 得第一项。

\textbf{用插值提高到次临界 $q$。}若 $p<q<p^*$，已证连续嵌入保证子列在 $L^{p^*}$ 中一致有界；选 $0<\theta<1$ 使 $1/q=\theta/p+(1-\theta)/p^*$。对差值插值，
\[
 \|u_n-u_m\|_q\le\|u_n-u_m\|_p^\theta
                    \|u_n-u_m\|_{p^*}^{1-\theta}\to0.
\]
若 $q<p$，直接用有限测度包含即可。临界 $p=d$ 时，对任意指定 $q>d$ 先选一个有限 $s>q$，用已证 $W^{1,d}\hookrightarrow L^s$，再在 $L^d$ 和 $L^s$ 间插值。

\textbf{超临界 H\"older 紧性。}Morrey 给出一致的 $C^{0,\alpha}$ 界，$\alpha=1-d/p$。在紧集 $\overline\Omega$ 上，用同样的有限网格方法得到一致收敛子列。对任意函数差 $w$，有
\[
 [w]_{C^{0,\beta}}
 \le C_{\alpha,\beta}\|w\|_\infty^{1-\beta/\alpha}
                           [w]_{C^{0,\alpha}}^{\beta/\alpha}.
\]
证明只需把 $|w(x)-w(y)|$ 同时控制为 $2\|w\|_\infty$ 与 $[w]_\alpha|x-y|^\alpha$，在两界相等的距离处分段。由于一致收敛而 $\alpha$ 半范数有界，便得 $C^{0,\beta}$ 强收敛。
\end{proof}

对任意有界开集的 $W^{1,p}_0$，同样的紧性结论可用零延拓证明，而不需要一般边界延拓。无界区域会因平移丧失紧性：取 $\phi(x-ne_1)$ 即可。临界指数处则有前面的集中反例。因此紧性需要同时防止向远处逃逸和向小尺度集中。

\subsection{零均值 Poincar\'e：用刚刚证明的紧性完成反证}

\begin{theorem}[零均值 Poincar\'e]\label{thm:sobolev-mean-poincare}
若 $\Omega$ 有界、连通且 $C^1$，$1\le p<\infty$，令 $u_\Omega=|\Omega|^{-1}\int_\Omega u$，则
\[
 \|u-u_\Omega\|_p\le C\|\nabla u\|_p.
\]
\end{theorem}
\begin{proof}
若不存在统一 $C$，则可取零均值 $v_n$，使 $\|v_n\|_p=1$、$\|\nabla v_n\|_p\to0$，方法是将反例减均值后按 $L^p$ 范数归一化。它们在 $W^{1,p}$ 中有界，由 Rellich 有子列在 $L^p$ 中强收敛于 $v$，故 $\|v\|_p=1$。另一方面测试函数恒等式取极限说明 $D_i v=0$，连通性与零梯度判据给出 $v$ 为常数。有限测度下 $L^p$ 强收敛推出积分收敛，所以 $\int v=0$，该常数只能为零，与范数为 1 矛盾。
\end{proof}
此证明明确按“零边界 Poincar\'e、独立证明 Rellich、再证明零均值 Poincar\'e”的顺序进行，没有循环依赖。

\subsection{高阶嵌入与乘积估计：以三维 \texorpdfstring{$H^2$}{H2} 为例}

高阶空间可通过逐次对导数应用一阶嵌入来理解。为避免额外引用高阶边界延拓，这里先在全空间 $\R^3$ 证明。若 $u\in H^2(\R^3)$，则 $u$ 和每个 $D_i u$ 均属于 $H^1$，因此次临界估计给出
\[
 \|u\|_6+\sum_i\|D_i u\|_6\le C\|u\|_{H^2}.
\]
于是 $u\in W^{1,6}$，再由 Morrey（$1-3/6=1/2$）得到
\begin{equation}
 H^2(\R^3)\hookrightarrow C^{0,1/2}(\R^3)\cap L^\infty(\R^3),
 \qquad\|u\|_\infty\le C\|u\|_{H^2}.
 \label{eq:sobolev-h2-embedding}
\end{equation}
同时 $\nabla u\in L^2\cap L^6$，插值得 $\|\nabla u\|_4\le C\|u\|_{H^2}$。对光滑 $u,v$，二阶乘积公式包含
\[
 D_{ij}(uv)=uD_{ij}v+vD_{ij}u+(D_i u)(D_j v)+(D_j u)(D_i v).
\]
前两项用 $L^\infty\cdot L^2$，后两项用 $L^4\cdot L^4$，连同零阶和一阶项得到
\begin{equation}
 \|uv\|_{H^2(\R^3)}\le C\|u\|_{H^2}\|v\|_{H^2}.
 \label{eq:sobolev-h2-algebra}
\end{equation}
通过全空间光滑逼近及该双线性估计可推广到一般 $H^2$ 函数。全空间的 $H^2$ 稠密性由前述截断与卷积对二阶导数重复即可证明。故三维 $H^2$ 是乘法代数，这对非线性强解估计很有用。

相比之下，三维 $H^1$ 一般不是乘法代数：取 $u=\chi(x)|x|^{-a}$，$1/4\le a<1/2$，$\chi$ 在原点附近为 1。前面的径向计算说明 $u\in H^1$，但 $u^2$ 的导数平方积分含 $\int_0^\delta r^{-4a}\dd r$，故 $u^2\notin H^1$。这明确说明何时不能直接把非线性乘积留在同一个空间。

\subsection{分数阶空间：与整数阶弱导数的关系}

对 $0<s<1$、$1\le p<\infty$，一种自然的分数阶定义是
\[
 [u]_{W^{s,p}(\Omega)}^p
 =\iint_{\Omega\times\Omega}\frac{|u(x)-u(y)|^p}{|x-y|^{d+sp}}\dd x\dd y,
 \quad\|u\|_{W^{s,p}}^p=\|u\|_p^p+[u]_{W^{s,p}}^p.
\]
它把差商在所有尺度上的平均大小作为“分数阶导数”，并非先要求经典非整数次求导。这里只给出定义和下面可直接验证的 Fourier 对应。

在 $\R^d$ 上采用酉归一化 Fourier 变换。对光滑快速衰减 $u$，令 $h=y-x$，由 Tonelli 与 Plancherel，
\begin{align*}
 [u]_{W^{s,2}(\R^d)}^2
 &=\int\frac{\|u(\cdot+h)-u\|_2^2}{|h|^{d+2s}}\dd h\\
 &=\int|\widehat u(\xi)|^2
       \left(\int\frac{|e^{i\xi\cdot h}-1|^2}{|h|^{d+2s}}\dd h\right)\dd\xi\\
 &=c_{d,s}\int|\xi|^{2s}|\widehat u(\xi)|^2\dd\xi.
\end{align*}
最后一步用旋转和缩放 $h\mapsto h/|\xi|$；常数有限，因为近原点分子为 $O(|h|^2)$，远处有界，径向积分分别是 $\int_0^1r^{1-2s}\dd r$ 和 $\int_1^\infty r^{-1-2s}\dd r$，恰好要求 $0<s<1$。

更一般定义 $H^s(\R^d)$ 的 Fourier 范数为 $\int(1+|\xi|^2)^s|\widehat u|^2$。在非负整数阶，它与弱导数的 $H^k$ 范数等价，因为 $\widehat{D^\alpha u}=(i\xi)^\alpha\widehat u$。这里使用 Fourier--Plancherel 理论作为另一路基础工具。分数阶空间解释了为何边界迹自然出现半阶正则性；一般区域上的分数阶迹和延拓不在这里证明。

\subsection{嵌入结论总表与时间紧性的边界}

\begin{center}\small
\begin{tabular}{p{0.18\linewidth}p{0.34\linewidth}p{0.35\linewidth}}
\toprule
条件 & 连续控制 & 紧性结论\\
\midrule
$1\le p<d$ & $W^{1,p}\hookrightarrow L^{p^*}$ & $L^q$ 紧，$1\le q<p^*$\\
$p=d\ge2$ & 每个有限 $L^q$ & 每个有限 $L^q$ 紧\\
$d<p<\infty$ & $C^{0,\alpha}$，$\alpha=1-d/p$ & $C^{0,\beta}$ 紧，$0<\beta<\alpha$\\
$d=1,p=1$ & 绝对连续且连续 & $L^q$ 紧，任意有限 $q$\\
\bottomrule
\end{tabular}
\end{center}
一般边界版本在本节均假设有界 $C^1$；零边界版本可采用零延拓。一维 $p=1$ 的有限 $L^q$ 紧性来自 $L^1$ 紧性、由绝对连续公式得到的统一 $L^\infty$ 界，以及插值；不宣称在 $C^0$ 中紧，因为跳变层可收缩。

空间紧性本身不能排除时间振荡，例如 $u_n(t,x)=\sin(nt)\phi(x)$。Aubin--Lions 型结论要求 $X_0\Subset X\hookrightarrow X_1$，$u_n$ 在 $L^a_tX_0$ 有界，且 $\partial_tu_n$ 在 $L^b_tX_1$ 有界（这里取 $1<a,b<\infty$），才得到 $L^a_tX$ 强收敛子列。该一般时间紧性定理在此明确作为外部定理引用；下一节会另证一个适用于常见三维 Galerkin 序列的 Hilbert 空间版本。

